\documentclass[12pt,a4paper]{article}
\usepackage[utf8]{vietnam}
\usepackage{amsmath,amssymb}
\usepackage{tabularx}
\usepackage{tikz}
\usepackage{xcolor}
\usepackage[top=0.6cm, bottom=0.6cm, left=0.6cm, right=0.6cm]{geometry}
\usepackage[dethi]{ex_test}

\SAOPTN{widthSA=2.8,heightSA=0.62}

\setlength{\parskip}{0pt}

% Dòng kẻ chấm mờ đều đặn chuẩn viết tay (khoảng cách 6.5mm rộng rãi vừa vặn)
\newcommand{\dongke}[1]{%
  \begingroup%
  \setlength{\parskip}{0pt}%
  \baselineskip=18.5pt%
  \vspace{1.0mm}%
  \foreach \i in {1,...,#1}{%
    \noindent{\color{gray!60}\dotfill}\par%
  }%
  \vspace{0.6mm}%
  \endgroup%
}

% Định nghĩa ô bài tập tự luận trong bảng
\newcommand{\oBaiTap}[3][6]{%
  \vspace{1.0mm}%
  \noindent\textbf{#2} #3\par%
  \dongke{#1}%
}

\begin{document}
\pagestyle{empty}
% =============================================================================
% ========================= TOÁN 6 - TUẦN 08 - BUỔI 01: MÃ ĐỀ 101 =============
% ========================= MẶT TRƯỚC (TRANG 1/2) ==============================
% =============================================================================

\noindent
\begin{tabularx}{\linewidth}{|X|c|}
\hline
\textbf{PHIẾU KIỂM TRA ĐỊNH KỲ TOÁN 6 - TUẦN 08 - BUỔI 01} & \textbf{Mã đề: 101 - Trang 1/2} \\
\hline
\end{tabularx}

\vspace{1mm}

\noindent
\begin{tabularx}{\linewidth}{|p{3.2cm}|X|}
\hline
\centering\textbf{ĐIỂM SỐ} & \centering\textbf{LỜI NHẬN XÉT CỦA THẦY CÔ} \tabularnewline
\hline
\rule{0pt}{1.8cm} & \tabularnewline
\hline
\end{tabularx}

\vspace{1.2mm}

\noindent\textbf{Câu 1 (0,25 điểm).} Tập hợp $A = \{x \in \mathbb{N} \mid 10 \le x \le 45\}$ có bao nhiêu phần tử? \shortans{36}\par\vspace{0.3mm}
\noindent\textbf{Câu 2 (0,25 điểm).} Số La Mã $\text{XXVIII}$ có giá trị trong hệ thập phân là bao nhiêu? \shortans{28}\par\vspace{0.3mm}
\noindent\textbf{Câu 3 (0,25 điểm).} Viết gọn tích $3^4 \cdot 3^3$ dưới dạng một lũy thừa. \shortans{3^7}\par\vspace{0.3mm}
\noindent\textbf{Câu 4 (0,25 điểm).} Giá trị của biểu thức $50 - 3 \cdot 2^3$ bằng bao nhiêu? \shortans{26}\par\vspace{0.3mm}
\noindent\textbf{Câu 5 (0,25 điểm).} Số dư của phép chia tổng $(35 + 120 + 43)$ cho $5$ là bao nhiêu? \shortans{3}\par\vspace{0.3mm}
\noindent\textbf{Câu 6 (0,25 điểm).} Tìm chữ số $x$ để số $\overline{45x}$ chia hết cho $9$ ($x > 0$). \shortans{9}\par\vspace{0.3mm}
\noindent\textbf{Câu 7 (0,25 điểm).} Trong các số $13; 21; 27; 35$, số nguyên tố là số nào? \shortans{13}\par\vspace{0.3mm}
\noindent\textbf{Câu 8 (0,25 điểm).} Một hình lục giác đều có bao nhiêu đường chéo chính? \shortans{3}\par\vspace{0.3mm}
\noindent\textbf{Câu 9 (0,25 điểm).} Một mảnh đất hình thoi có độ dài cạnh là $12\text{ m}$. Chu vi là: \shortans{48\text{ m}}\par\vspace{0.3mm}
\noindent\textbf{Câu 10 (0,25 điểm).} Mảnh bìa hình thoi có hai đường chéo là $10\text{ cm}$ và $14\text{ cm}$. Diện tích là: \shortans{70\text{ cm}^2}\par\vspace{0.3mm}
\vspace{0.4mm}

\noindent\textbf{Câu 11 (3,0 điểm). Thực hiện phép tính (hợp lý nếu có thể):}
\vspace{0.6mm}

\noindent
\begin{tabularx}{\linewidth}{|X|X|}
\hline
\oBaiTap[7]{a)}{145 - 2^3 \cdot 5 + 3^2 \cdot 4} & 
\oBaiTap[7]{b)}{38 + 147 + 62 + 53} \\
\hline
\oBaiTap[7]{c)}{48 \cdot 67 + 48 \cdot 33 - 800} & 
\oBaiTap[7]{d)}{180 : (25 - 7) + 3^2 \cdot 5} \\
\hline
\oBaiTap[7]{e)}{180 : [45 - (18 - 15)^2]} & 
\oBaiTap[7]{f)}{250 : \{ 90 - [32 + (14 - 12)^3] \}} \\
\hline
\end{tabularx}

\newpage
% =============================================================================
% ========================= TOÁN 6 - TUẦN 08 - BUỔI 01: MÃ ĐỀ 101 =============
% ========================= MẶT SAU (TRANG 2/2) ================================
% =============================================================================

\noindent
\begin{tabularx}{\linewidth}{|X|c|}
\hline
\textbf{TOÁN 6 - TUẦN 08 - BUỔI 01} & \textbf{Mã đề: 101 - Trang 2/2} \\
\hline
\end{tabularx}

\vspace{1.5mm}

\noindent\textbf{Câu 12 (2,0 điểm). Tìm số tự nhiên $x$, biết:}
\vspace{0.8mm}

\noindent
\begin{tabularx}{\linewidth}{|X|X|}
\hline
\oBaiTap[6]{a)}{4 \cdot x + 15 = 47} & 
\oBaiTap[6]{b)}{85 - 3 \cdot (x + 4) = 40} \\
\hline
\oBaiTap[6]{c)}{3 \cdot 2^{x+1} - 12 = 36} & 
\oBaiTap[6]{d)}{\text{Tìm } x \text{ lớn nhất biết } 24 \ \vdots \ x \text{ và } 36 \ \vdots \ x} \\
\hline
\end{tabularx}

\vspace{1.5mm}

\noindent
\begin{tabularx}{\linewidth}{|X|}
\hline
\vspace{0.5mm}
\textbf{Câu 13 (1,5 điểm).} Bác Ba có một mảnh vườn trồng rau hình chữ nhật có chiều dài $15\text{ m}$ và chiều rộng $8\text{ m}$. Bác Ba dự định làm hàng rào lưới thép bao quanh khu vườn (chừa lại cổng ra vào rộng $2\text{ m}$).\par
\textbf{a) (0,75 điểm)} Tính chu vi và diện tích của mảnh vườn.\par
\textbf{b) (0,75 điểm)} Biết chi phí làm mỗi mét hàng rào là $120.000$ đồng. Hỏi bác Ba phải trả bao nhiêu tiền làm hàng rào?\par
\vspace{0.4mm}
\noindent\textit{Bài giải:}\dongke{11}
\\ \hline
\end{tabularx}

\vspace{1.5mm}

\noindent
\begin{tabularx}{\linewidth}{|X|}
\hline
\vspace{0.5mm}
\textbf{Câu 14 (1,0 điểm).} Cho tổng $A = 1 + 3 + 3^2 + 3^3 + \dots + 3^{21}$. Chứng minh rằng $A$ chia hết cho $4$.\par
\vspace{0.4mm}
\noindent\textit{Bài giải:}\dongke{8}
\\ \hline
\end{tabularx}

\newpage

% =============================================================================
% ========================= TOÁN 6 - TUẦN 08 - BUỔI 01: MÃ ĐỀ 102 =============
% ========================= MẶT TRƯỚC (TRANG 1/2) ==============================
% =============================================================================

\noindent
\begin{tabularx}{\linewidth}{|X|c|}
\hline
\textbf{PHIẾU KIỂM TRA ĐỊNH KỲ TOÁN 6 - TUẦN 08 - BUỔI 01} & \textbf{Mã đề: 102 - Trang 1/2} \\
\hline
\end{tabularx}

\vspace{1mm}

\noindent
\begin{tabularx}{\linewidth}{|p{3.2cm}|X|}
\hline
\centering\textbf{ĐIỂM SỐ} & \centering\textbf{LỜI NHẬN XÉT CỦA THẦY CÔ} \tabularnewline
\hline
\rule{0pt}{1.8cm} & \tabularnewline
\hline
\end{tabularx}

\vspace{1.2mm}

\noindent\textbf{Câu 1 (0,25 điểm).} Tập hợp $A = \{x \in \mathbb{N} \mid 12 \le x \le 55\}$ có bao nhiêu phần tử? \shortans{44}\par\vspace{0.3mm}
\noindent\textbf{Câu 2 (0,25 điểm).} Số La Mã $\text{XXIV}$ có giá trị trong hệ thập phân là bao nhiêu? \shortans{24}\par\vspace{0.3mm}
\noindent\textbf{Câu 3 (0,25 điểm).} Viết gọn tích $2^5 \cdot 2^4$ dưới dạng một lũy thừa. \shortans{2^9}\par\vspace{0.3mm}
\noindent\textbf{Câu 4 (0,25 điểm).} Giá trị của biểu thức $60 - 4 \cdot 3^2$ bằng bao nhiêu? \shortans{24}\par\vspace{0.3mm}
\noindent\textbf{Câu 5 (0,25 điểm).} Số dư của phép chia tổng $(45 + 130 + 54)$ cho $5$ là bao nhiêu? \shortans{4}\par\vspace{0.3mm}
\noindent\textbf{Câu 6 (0,25 điểm).} Tìm chữ số $x$ lớn nhất để số $\overline{32x}$ chia hết cho $3$. \shortans{7}\par\vspace{0.3mm}
\noindent\textbf{Câu 7 (0,25 điểm).} Trong các số $17; 25; 33; 49$, số nguyên tố là số nào? \shortans{17}\par\vspace{0.3mm}
\noindent\textbf{Câu 8 (0,25 điểm).} Một hình lục giác đều có bao nhiêu cạnh bằng nhau? \shortans{6}\par\vspace{0.3mm}
\noindent\textbf{Câu 9 (0,25 điểm).} Một mảnh đất hình thoi có độ dài cạnh là $15\text{ m}$. Chu vi là: \shortans{60\text{ m}}\par\vspace{0.3mm}
\noindent\textbf{Câu 10 (0,25 điểm).} Mảnh bìa hình thoi có hai đường chéo là $12\text{ cm}$ và $15\text{ cm}$. Diện tích là: \shortans{90\text{ cm}^2}\par\vspace{0.3mm}
\vspace{0.4mm}

\noindent\textbf{Câu 11 (3,0 điểm). Thực hiện phép tính (hợp lý nếu có thể):}
\vspace{0.6mm}

\noindent
\begin{tabularx}{\linewidth}{|X|X|}
\hline
\oBaiTap[7]{a)}{155 - 2^4 \cdot 3 + 4^2 \cdot 2} & 
\oBaiTap[7]{b)}{46 + 138 + 54 + 62} \\
\hline
\oBaiTap[7]{c)}{56 \cdot 72 + 56 \cdot 28 - 600} & 
\oBaiTap[7]{d)}{160 : (24 - 8) + 4^2 \cdot 3} \\
\hline
\oBaiTap[7]{e)}{210 : [58 - (19 - 15)^2]} & 
\oBaiTap[7]{f)}{300 : \{ 100 - [32 + (15 - 13)^3] \}} \\
\hline
\end{tabularx}

\newpage
% =============================================================================
% ========================= TOÁN 6 - TUẦN 08 - BUỔI 01: MÃ ĐỀ 102 =============
% ========================= MẶT SAU (TRANG 2/2) ================================
% =============================================================================

\noindent
\begin{tabularx}{\linewidth}{|X|c|}
\hline
\textbf{TOÁN 6 - TUẦN 08 - BUỔI 01} & \textbf{Mã đề: 102 - Trang 2/2} \\
\hline
\end{tabularx}

\vspace{1.5mm}

\noindent\textbf{Câu 12 (2,0 điểm). Tìm số tự nhiên $x$, biết:}
\vspace{0.8mm}

\noindent
\begin{tabularx}{\linewidth}{|X|X|}
\hline
\oBaiTap[6]{a)}{5 \cdot x + 18 = 53} & 
\oBaiTap[6]{b)}{95 - 4 \cdot (x + 3) = 47} \\
\hline
\oBaiTap[6]{c)}{2 \cdot 3^{x+1} - 15 = 39} & 
\oBaiTap[6]{d)}{\text{Tìm } x \text{ lớn nhất biết } 30 \ \vdots \ x \text{ và } 45 \ \vdots \ x} \\
\hline
\end{tabularx}

\vspace{1.5mm}

\noindent
\begin{tabularx}{\linewidth}{|X|}
\hline
\vspace{0.5mm}
\textbf{Câu 13 (1,5 điểm).} Bác Ba có một mảnh vườn trồng rau hình chữ nhật có chiều dài $16\text{ m}$ và chiều rộng $9\text{ m}$. Bác Ba dự định làm hàng rào lưới thép bao quanh khu vườn (chừa lại cổng ra vào rộng $2\text{ m}$).\par
\textbf{a) (0,75 điểm)} Tính chu vi và diện tích của mảnh vườn.\par
\textbf{b) (0,75 điểm)} Biết chi phí làm mỗi mét hàng rào là $150.000$ đồng. Hỏi bác Ba phải trả bao nhiêu tiền làm hàng rào?\par
\vspace{0.4mm}
\noindent\textit{Bài giải:}\dongke{11}
\\ \hline
\end{tabularx}

\vspace{1.5mm}

\noindent
\begin{tabularx}{\linewidth}{|X|}
\hline
\vspace{0.5mm}
\textbf{Câu 14 (1,0 điểm).} Cho tổng $B = 1 + 4 + 4^2 + 4^3 + \dots + 4^{21}$. Chứng minh rằng $B$ chia hết cho $5$.\par
\vspace{0.4mm}
\noindent\textit{Bài giải:}\dongke{8}
\\ \hline
\end{tabularx}

\newpage

% =============================================================================
% ========================= TOÁN 6 - TUẦN 08 - BUỔI 01: MÃ ĐỀ 103 =============
% ========================= MẶT TRƯỚC (TRANG 1/2) ==============================
% =============================================================================

\noindent
\begin{tabularx}{\linewidth}{|X|c|}
\hline
\textbf{PHIẾU KIỂM TRA ĐỊNH KỲ TOÁN 6 - TUẦN 08 - BUỔI 01} & \textbf{Mã đề: 103 - Trang 1/2} \\
\hline
\end{tabularx}

\vspace{1mm}

\noindent
\begin{tabularx}{\linewidth}{|p{3.2cm}|X|}
\hline
\centering\textbf{ĐIỂM SỐ} & \centering\textbf{LỜI NHẬN XÉT CỦA THẦY CÔ} \tabularnewline
\hline
\rule{0pt}{1.8cm} & \tabularnewline
\hline
\end{tabularx}

\vspace{1.2mm}

\noindent\textbf{Câu 1 (0,25 điểm).} Tập hợp $A = \{x \in \mathbb{N} \mid 15 \le x \le 63\}$ có bao nhiêu phần tử? \shortans{49}\par\vspace{0.3mm}
\noindent\textbf{Câu 2 (0,25 điểm).} Số La Mã $\text{XXIX}$ có giá trị trong hệ thập phân là bao nhiêu? \shortans{29}\par\vspace{0.3mm}
\noindent\textbf{Câu 3 (0,25 điểm).} Viết gọn tích $5^3 \cdot 5^4$ dưới dạng một lũy thừa. \shortans{5^7}\par\vspace{0.3mm}
\noindent\textbf{Câu 4 (0,25 điểm).} Giá trị của biểu thức $70 - 2 \cdot 4^2$ bằng bao nhiêu? \shortans{38}\par\vspace{0.3mm}
\noindent\textbf{Câu 5 (0,25 điểm).} Số dư của phép chia tổng $(55 + 140 + 62)$ cho $5$ là bao nhiêu? \shortans{2}\par\vspace{0.3mm}
\noindent\textbf{Câu 6 (0,25 điểm).} Tìm chữ số $x$ lớn nhất để số $\overline{51x}$ chia hết cho $3$. \shortans{9}\par\vspace{0.3mm}
\noindent\textbf{Câu 7 (0,25 điểm).} Trong các số $19; 39; 51; 57$, số nguyên tố là số nào? \shortans{19}\par\vspace{0.3mm}
\noindent\textbf{Câu 8 (0,25 điểm).} Một hình lục giác đều có bao nhiêu góc ở các đỉnh bằng nhau? \shortans{6}\par\vspace{0.3mm}
\noindent\textbf{Câu 9 (0,25 điểm).} Một mảnh đất hình thoi có độ dài cạnh là $14\text{ m}$. Chu vi là: \shortans{56\text{ m}}\par\vspace{0.3mm}
\noindent\textbf{Câu 10 (0,25 điểm).} Mảnh bìa hình thoi có hai đường chéo là $8\text{ cm}$ và $16\text{ cm}$. Diện tích là: \shortans{64\text{ cm}^2}\par\vspace{0.3mm}
\vspace{0.4mm}

\noindent\textbf{Câu 11 (3,0 điểm). Thực hiện phép tính (hợp lý nếu có thể):}
\vspace{0.6mm}

\noindent
\begin{tabularx}{\linewidth}{|X|X|}
\hline
\oBaiTap[7]{a)}{165 - 3^3 \cdot 2 + 5^2 \cdot 3} & 
\oBaiTap[7]{b)}{29 + 154 + 71 + 46} \\
\hline
\oBaiTap[7]{c)}{64 \cdot 75 + 64 \cdot 25 - 400} & 
\oBaiTap[7]{d)}{210 : (27 - 6) + 5^2 \cdot 2} \\
\hline
\oBaiTap[7]{e)}{220 : [69 - (21 - 16)^2]} & 
\oBaiTap[7]{f)}{350 : \{ 110 - [32 + (16 - 14)^3] \}} \\
\hline
\end{tabularx}

\newpage
% =============================================================================
% ========================= TOÁN 6 - TUẦN 08 - BUỔI 01: MÃ ĐỀ 103 =============
% ========================= MẶT SAU (TRANG 2/2) ================================
% =============================================================================

\noindent
\begin{tabularx}{\linewidth}{|X|c|}
\hline
\textbf{TOÁN 6 - TUẦN 08 - BUỔI 01} & \textbf{Mã đề: 103 - Trang 2/2} \\
\hline
\end{tabularx}

\vspace{1.5mm}

\noindent\textbf{Câu 12 (2,0 điểm). Tìm số tự nhiên $x$, biết:}
\vspace{0.8mm}

\noindent
\begin{tabularx}{\linewidth}{|X|X|}
\hline
\oBaiTap[6]{a)}{6 \cdot x + 14 = 56} & 
\oBaiTap[6]{b)}{90 - 5 \cdot (x + 2) = 45} \\
\hline
\oBaiTap[6]{c)}{4 \cdot 2^{x+2} - 18 = 46} & 
\oBaiTap[6]{d)}{\text{Tìm } x \text{ lớn nhất biết } 28 \ \vdots \ x \text{ và } 42 \ \vdots \ x} \\
\hline
\end{tabularx}

\vspace{1.5mm}

\noindent
\begin{tabularx}{\linewidth}{|X|}
\hline
\vspace{0.5mm}
\textbf{Câu 13 (1,5 điểm).} Bác Ba có một mảnh vườn trồng rau hình chữ nhật có chiều dài $18\text{ m}$ và chiều rộng $10\text{ m}$. Bác Ba dự định làm hàng rào lưới thép bao quanh khu vườn (chừa lại cổng ra vào rộng $2\text{ m}$).\par
\textbf{a) (0,75 điểm)} Tính chu vi và diện tích của mảnh vườn.\par
\textbf{b) (0,75 điểm)} Biết chi phí làm mỗi mét hàng rào là $140.000$ đồng. Hỏi bác Ba phải trả bao nhiêu tiền làm hàng rào?\par
\vspace{0.4mm}
\noindent\textit{Bài giải:}\dongke{11}
\\ \hline
\end{tabularx}

\vspace{1.5mm}

\noindent
\begin{tabularx}{\linewidth}{|X|}
\hline
\vspace{0.5mm}
\textbf{Câu 14 (1,0 điểm).} Cho tổng $C = 1 + 5 + 5^2 + 5^3 + \dots + 5^{21}$. Chứng minh rằng $C$ chia hết cho $6$.\par
\vspace{0.4mm}
\noindent\textit{Bài giải:}\dongke{8}
\\ \hline
\end{tabularx}

\newpage

% =============================================================================
% ========================= TOÁN 6 - TUẦN 08 - BUỔI 01: MÃ ĐỀ 104 =============
% ========================= MẶT TRƯỚC (TRANG 1/2) ==============================
% =============================================================================

\noindent
\begin{tabularx}{\linewidth}{|X|c|}
\hline
\textbf{PHIẾU KIỂM TRA ĐỊNH KỲ TOÁN 6 - TUẦN 08 - BUỔI 01} & \textbf{Mã đề: 104 - Trang 1/2} \\
\hline
\end{tabularx}

\vspace{1mm}

\noindent
\begin{tabularx}{\linewidth}{|p{3.2cm}|X|}
\hline
\centering\textbf{ĐIỂM SỐ} & \centering\textbf{LỜI NHẬN XÉT CỦA THẦY CÔ} \tabularnewline
\hline
\rule{0pt}{1.8cm} & \tabularnewline
\hline
\end{tabularx}

\vspace{1.2mm}

\noindent\textbf{Câu 1 (0,25 điểm).} Tập hợp $A = \{x \in \mathbb{N} \mid 18 \le x \le 72\}$ có bao nhiêu phần tử? \shortans{55}\par\vspace{0.3mm}
\noindent\textbf{Câu 2 (0,25 điểm).} Số La Mã $\text{XXVII}$ có giá trị trong hệ thập phân là bao nhiêu? \shortans{27}\par\vspace{0.3mm}
\noindent\textbf{Câu 3 (0,25 điểm).} Viết gọn tích $7^2 \cdot 7^5$ dưới dạng một lũy thừa. \shortans{7^7}\par\vspace{0.3mm}
\noindent\textbf{Câu 4 (0,25 điểm).} Giá trị của biểu thức $80 - 5 \cdot 2^3$ bằng bao nhiêu? \shortans{40}\par\vspace{0.3mm}
\noindent\textbf{Câu 5 (0,25 điểm).} Số dư của phép chia tổng $(65 + 150 + 71)$ cho $5$ là bao nhiêu? \shortans{1}\par\vspace{0.3mm}
\noindent\textbf{Câu 6 (0,25 điểm).} Tìm chữ số $x$ để số $\overline{63x}$ chia hết cho $9$ ($x > 0$). \shortans{9}\par\vspace{0.3mm}
\noindent\textbf{Câu 7 (0,25 điểm).} Trong các số $23; 35; 45; 63$, số nguyên tố là số nào? \shortans{23}\par\vspace{0.3mm}
\noindent\textbf{Câu 8 (0,25 điểm).} Một hình vuông có bao nhiêu trục đối xứng? \shortans{4}\par\vspace{0.3mm}
\noindent\textbf{Câu 9 (0,25 điểm).} Một mảnh đất hình thoi có độ dài cạnh là $16\text{ m}$. Chu vi là: \shortans{64\text{ m}}\par\vspace{0.3mm}
\noindent\textbf{Câu 10 (0,25 điểm).} Mảnh bìa hình thoi có hai đường chéo là $14\text{ cm}$ và $18\text{ cm}$. Diện tích là: \shortans{126\text{ cm}^2}\par\vspace{0.3mm}
\vspace{0.4mm}

\noindent\textbf{Câu 11 (3,0 điểm). Thực hiện phép tính (hợp lý nếu có thể):}
\vspace{0.6mm}

\noindent
\begin{tabularx}{\linewidth}{|X|X|}
\hline
\oBaiTap[7]{a)}{175 - 2^5 \cdot 2 + 3^3 \cdot 3} & 
\oBaiTap[7]{b)}{37 + 165 + 63 + 35} \\
\hline
\oBaiTap[7]{c)}{72 \cdot 81 + 72 \cdot 19 - 200} & 
\oBaiTap[7]{d)}{240 : (29 - 5) + 6^2 \cdot 2} \\
\hline
\oBaiTap[7]{e)}{240 : [84 - (23 - 17)^2]} & 
\oBaiTap[7]{f)}{400 : \{ 120 - [32 + (17 - 15)^3] \}} \\
\hline
\end{tabularx}

\newpage
% =============================================================================
% ========================= TOÁN 6 - TUẦN 08 - BUỔI 01: MÃ ĐỀ 104 =============
% ========================= MẶT SAU (TRANG 2/2) ================================
% =============================================================================

\noindent
\begin{tabularx}{\linewidth}{|X|c|}
\hline
\textbf{TOÁN 6 - TUẦN 08 - BUỔI 01} & \textbf{Mã đề: 104 - Trang 2/2} \\
\hline
\end{tabularx}

\vspace{1.5mm}

\noindent\textbf{Câu 12 (2,0 điểm). Tìm số tự nhiên $x$, biết:}
\vspace{0.8mm}

\noindent
\begin{tabularx}{\linewidth}{|X|X|}
\hline
\oBaiTap[6]{a)}{7 \cdot x + 16 = 65} & 
\oBaiTap[6]{b)}{100 - 6 \cdot (x + 1) = 52} \\
\hline
\oBaiTap[6]{c)}{5 \cdot 3^{x-1} + 15 = 60} & 
\oBaiTap[6]{d)}{\text{Tìm } x \text{ lớn nhất biết } 32 \ \vdots \ x \text{ và } 48 \ \vdots \ x} \\
\hline
\end{tabularx}

\vspace{1.5mm}

\noindent
\begin{tabularx}{\linewidth}{|X|}
\hline
\vspace{0.5mm}
\textbf{Câu 13 (1,5 điểm).} Bác Ba có một mảnh vườn trồng rau hình chữ nhật có chiều dài $20\text{ m}$ và chiều rộng $11\text{ m}$. Bác Ba dự định làm hàng rào lưới thép bao quanh khu vườn (chừa lại cổng ra vào rộng $2\text{ m}$).\par
\textbf{a) (0,75 điểm)} Tính chu vi và diện tích của mảnh vườn.\par
\textbf{b) (0,75 điểm)} Biết chi phí làm mỗi mét hàng rào là $130.000$ đồng. Hỏi bác Ba phải trả bao nhiêu tiền làm hàng rào?\par
\vspace{0.4mm}
\noindent\textit{Bài giải:}\dongke{11}
\\ \hline
\end{tabularx}

\vspace{1.5mm}

\noindent
\begin{tabularx}{\linewidth}{|X|}
\hline
\vspace{0.5mm}
\textbf{Câu 14 (1,0 điểm).} Cho tổng $D = 1 + 2 + 2^2 + 2^3 + \dots + 2^{23}$. Chứng minh rằng $D$ chia hết cho $3$.\par
\vspace{0.4mm}
\noindent\textit{Bài giải:}\dongke{8}
\\ \hline
\end{tabularx}

\newpage

\end{document}